% Table 11: amplitude output, forward-model residual and phase-scale probes.
\begin{table}[!htbp]
\centering
\caption{(a) Amplitude output against the reconstruction-derived amplitude
reference (113 off-axis test fields; one run each). (b) Forward-model residual
on the 113 off-axis test fields (Baseline and +IPP (per-cell): mean $\pm$ SD
over three seeds; others one run). (c) Change in the residual when the
reference phase is scaled by 0.9 or 0.5 (margin: scaled minus reference;
positive means the reference phase has the lower residual).}
\label{tab:amplitude}
\footnotesize
\textbf{(a) Amplitude output}\\[2pt]
\begin{tabular*}{\textwidth}{@{\extracolsep{\fill}}lrrr@{}}
\toprule
Quantity & +Amplitude & +Fwd (fixed $z$) & +Fwd (free $z$) \\
\midrule
MAE vs the reference & 0.0651 & 0.0665 & 0.0651 \\
\quad the same MAE for $A=1$ & 0.1469 & 0.1469 & 0.1469 \\
\quad ratio ($<1$: closer than $A=1$) & 0.4429 & 0.4526 & 0.4436 \\
MAE inside cells & 0.0906 & 0.0922 & 0.0898 \\
Bias & $+0.0339$ & $+0.0378$ & $+0.0323$ \\
Pearson $r$ & 0.9047 & 0.9064 & 0.9045 \\
Predicted amplitude, mean & 1.0324 & 1.0363 & 1.0308 \\
Predicted amplitude, SD & 0.1207 & 0.1223 & 0.1200 \\
Reference amplitude, mean & 0.9985 & 0.9985 & 0.9985 \\
\bottomrule
\end{tabular*}

\vspace{8pt}
\textbf{(b) Forward-model residual}\\[2pt]
\begin{tabular*}{\textwidth}{@{\extracolsep{\fill}}lccccc@{}}
\toprule
Quantity & Baseline & +IPP (per-cell) & +Amplitude & +Fwd (fixed $z$) & +Fwd (free $z$) \\
\midrule
Residual, prediction & $0.8948 \pm 0.0002$ & $0.8968 \pm 0.0005$ & 0.8699 & 0.8702 & 0.8699 \\
Residual, reference phase & $0.8914 \pm 0.0000$ & $0.8914 \pm 0.0000$ & 0.8647 & 0.8643 & 0.8640 \\
Ratio (1.0 = reference phase) & $1.0039 \pm 0.0002$ & $1.0061 \pm 0.0005$ & 1.0061 & 1.0068 & 1.0068 \\
In-cell phase MAE [rad] & $0.2689 \pm 0.0056$ & $0.2876 \pm 0.0064$ & 0.2862 & 0.2893 & 0.2886 \\
\bottomrule
\end{tabular*}

\vspace{8pt}
\textbf{(c) Phase-scale probes}\\[2pt]
\begin{tabular*}{\textwidth}{@{\extracolsep{\fill}}l>{\raggedleft\arraybackslash}p{1.5cm}rp{2.4cm}rp{2.4cm}@{}}
\toprule
 & & \multicolumn{2}{c}{$0.9\times$ phase} & \multicolumn{2}{c}{$0.5\times$ phase} \\
\cmidrule(lr){3-4}\cmidrule(lr){5-6}
Geometry and fields & Residual & Margin & Fields favouring reference phase & Margin & Fields favouring reference phase \\
\midrule
Off-axis, 32 validation fields & 0.9147 & $+0.00035$ & 18/32 & $+0.0080$ & 20/32 \\
In-line, 32 validation fields & 0.6993 & $+0.00432$ & 30/32 & $+0.0622$ & 32/32 \\
Off-axis, 113 test fields, global surface & 0.8914 & $+0.00027$ & 60/113 & $+0.0070$ & 71/113 \\
Off-axis, 112 test fields, per-field surfaces$^{\ast}$ & 0.3906 & $+0.00110$ & 93/112 & $+0.0359$ & 110/112 \\
\bottomrule
\end{tabular*}
\tabnote{Fields favouring reference phase: number of fields on which the scaled phase gives the higher residual than the reference
phase (a count of fields, not a separate test). Margin: scaled minus reference residual, averaged over the fields. Configured tolerance for
a usable margin: 0.01. Residual levels in (b) for configurations with an amplitude output use the predicted amplitude in place of
$A=1$ and are not comparable with the others; only the ratio is. $^{\ast}$Per-field aberration surfaces are fitted against each
field's reference phase and are not available at inference. One test field (T24\_Staurosporine\_100nM\_10) has
no per-field surface: its fit was rejected because the fitted surface exceeded the 45\,rad peak-to-valley limit (405\,rad).
With the global surface on the same 112 fields the residual is 0.8926, the margins are
$+0.00025$ and $+0.0067$ and the counts 59/112 and 70/112.}
\end{table}
